\section{Related Work}\label{sec:related}

\stitle{Enumeration-based core and nucleus decomposition.}
The \(k\)-core was introduced by Seidman~\cite{kcoreSeidman}.
%
Batagelj and Zaver\v{s}nik~\cite{coreAlg} gave the linear-time peeling algorithm with a bucket queue.
%
Nucleus decomposition generalizes this idea from edges to higher-order clique witnesses~\cite{nucleusSariyuce14}.
%
In the general \((r,s)\) setting, peeling removes \(r\)-cliques according to their support in \(s\)-cliques.
%
That generality is broader than the problem solved here.
%
Our work specializes to the vertex-centric \(r=1\) case and exploits a state invariant that does not hold for general \(r\).

\stitle{Compact clique representations.}
Bron--Kerbosch search~\cite{BronKerbosch}, Tomita pivoting~\cite{tomita}, and degeneracy ordering~\cite{eppstein} are standard tools for clique enumeration.
%
Pivoter introduced the succinct clique tree as a compact representation for clique counting~\cite{Pivoter}.
%
The Clique Path Index builds on this line to support general nucleus decomposition without explicit \(s\)-clique enumeration~\cite{NuclearCD}.
%
Those representations answer how to compactly encode many cliques.
%
This paper asks a different question: once the compact index exists, does \(r=1\) require it to remain mutable during peeling?
%
The static-\cpi{} theorem answers no.

\stitle{Local and parallel peeling.}
\(h\)-index characterizations of core values were studied for local core computation~\cite{lu2016hcore,kcoreLocal,hindex}.
%
\localH{} is the direct lift of that fixed-point view to the implicit \(s\)-clique hypergraph encoded by the \cpi{}.
%
\peelH{} is the peel-order realization of the same equations.
%
The difference from prior local \(h\)-index work is not the operator itself.
%
The difference is that the witness multiset is represented by static \cpi{} paths and updated by binomial counter deltas.

\stitle{Parallel construction and clique counting.}
Parallel clique search usually tries to distribute the expensive recursive enumeration work.
%
\paraBuild{} follows the same seed-level decomposition principle.
%
Its specific role in this paper is tied to the static consumer layout.
%
Because the downstream peel never mutates path sets, construction can stream paths into thread-local arenas and merge flat incidence arrays.
%
This is an engineering consequence of the \(r=1\) state model rather than a new clique-enumeration paradigm.

\stitle{Hierarchical cohesive subgraphs.}
Nucleus decompositions naturally form hierarchies across core levels~\cite{nucleusSariyuce14}.
%
Elder-rule join-tree extraction is standard in computational topology~\cite{morozov2007persistence}.
%
Prior nucleus pipelines can extract such hierarchies after sorting decomposition output.
%
\buildHier{} avoids the extra sort because \peelH{} already emits a nondecreasing deletion log.
%
The hierarchy result is therefore a byproduct of the chosen peel realization.

\stitle{Empirical graph mining methodology.}
The evaluation follows common practice in cohesive-subgraph experiments: SNAP community ground truth, large public graphs, runtime and memory comparisons, and dense synthetic stress tests~\cite{communityFinder,NuclearCD,nucleusSariyuceTODS}.
%
The case studies use those benchmarks to test whether the \(r=1\) specialization remains useful as a vertex-ranking tool.
